% !TeX program = tectonic
\documentclass[11pt,a4paper]{letter}
\usepackage[margin=1in]{geometry}
\usepackage{enumitem}
\setlist{nosep}
\setlength{\parskip}{0pt}

\signature{[Corresponding author name]}

\begin{document}

\begin{letter}{Editorial Office\\Nature Communications\\Springer Nature}

\opening{Dear Editors,}
\vspace{-1em}

\noindent\textbf{Submission:} Profile-Driven Hardware Simulation for Organic Optoelectronic Edge Vision\\
\textbf{Article Type:} Article

\textbf{Editorial Summary}

We present an open-source simulation framework for organic optoelectronic compute-in-memory inference with modern vision transformers. Our central contribution, Ensemble HAT, resamples device-to-device mismatch masks every epoch. On Tiny-ViT-5M, it raises fresh-instance accuracy from chance level (10.00\%) to 86.37\% (mean ± 1.54\%), enabling robust deployment without per-device calibration.

We rigorously falsify three intuitive strategies for breaking the severe nonlinearity ceiling---joint linearization training, all-linear ablations, and MLP-only controls---all converging to an approximately 30\% fresh-instance limit. Rather than weakening the study, this negative result exposes a structural generalization barrier in the attention pathway that first-order behavioral surrogates cannot circumvent.

\textbf{Novelty Claims}
\begin{itemize}
    \item We introduce the first open-source framework for organic optoelectronic CIM with vision transformers, providing PyTorch-native device-profile substitution and hybrid analog/digital deployment.
    \item We identify hardware-instance overfitting in organic CIM inference and propose Ensemble HAT, a training-time mismatch-mask resampling fix that restores generalization.
    \item We falsify three severe-nonlinearity mitigation strategies, establishing an approximately 30\% structural ceiling that defines the boundary of first-order behavioral simulation.
\end{itemize}

\textbf{Audience}

Nature Communications readers will find a reusable methodology for evaluating organic device assumptions on vision tasks and a rigorous negative-result anchor that delineates the limits of first-order behavioral simulation.

\textbf{Related Submissions}

This work is not under consideration elsewhere.

\textbf{Suggested Reviewers}

[To be provided]

\textbf{Contact Information}

Li Songqiao / 2622507532@qq.com

\closing{Sincerely,}

\end{letter}
\end{document}
